\chapter{Future Directions}
\label{ch:future}

The world-model calculus establishes a foundation.  Several directions
extend it in ways that are natural continuations of the existing
development rather than speculative leaps.

\section{Measure-Theoretic Grounding}
\label{sec:future-measure}

The current formalization works over finite world-sets.
Measure-theoretic extensions would ground the calculus on
$\sigma$-algebras and probability measures, enabling continuous-valued
evidence, non-conjugate distributional families, and integration with
the broader mathematical probability literature.  The kernel's
compositionality law generalizes naturally to the measure-theoretic
setting; the technical challenge is in the evidence-carrier
instantiation and the convergence bridges.

\section{HOL Completeness}

The natively typed world-model layer
(Chapter~\ref{ch:typed}) establishes syntax, soundness, and type
synthesis.  A typed canonical-model completeness theorem would close
the gap between the calculus's higher-order syntax and full semantic
adequacy.  The required construction---a term model over equivalence
classes of typed evidence---is well understood in principle but
technically demanding.

\section{Selective Distributional Inference}

The convergence and dominance results in the error-transport chapter
(Chapter~\ref{ch:error-transport}) show that scalar inference is
asymptotically correct but suboptimal in the finite-evidence regime.
The natural next step is a cost-aware inference controller that
escalates from scalar to distributional propagation when the regime
structure predicts a significant gain---using compact approximation
families, compiled rule surrogates, and accelerator-backed batched
computation.

\section{Non-Additive Regimes}

The evidence algebra's additive structure
(Chapter~\ref{ch:evidence}) is appropriate for independent
observations but does not capture all forms of evidential
combination.  Non-additive regimes---belief functions,
possibility measures, imprecise probabilities beyond the Walley
framework---represent a natural generalization.  The challenge is to
preserve compositionality while relaxing additivity, and the
quantale structure suggests a path: replace the additive monoid with
a more general tensor.

\section{Probability Hypercube}

The current view projections extract one-dimensional summaries
(strength, confidence) from the evidence space.  A probability
hypercube representation would retain the full joint structure of
multi-query evidence, enabling richer decision-making and more
informative uncertainty quantification.  This connects to the
infinitary extension: the hypercube is naturally a product measure
space.

\section{Universality}

The typed world-model encoding
(Chapter~\ref{ch:typed}) realizes the calculus as a specific
language definition.  A universality result would show that any typed
language definition satisfying certain compositionality axioms admits
a world-model interpretation---establishing the calculus as a
canonical representative of a broader class rather than one design
among many.

\section{Operational Runtime}

The formalization specifies the kernel semantics precisely, but a
high-performance compiled runtime---executing the same
posterior-state calculus at scale with caching, incremental evidence
fusion, and multi-source integration---is the bridge from formalized
semantics to deployed reasoning.  The $\Sigma$-guarded rewrite
system (Chapter~\ref{ch:sigma-rewrites}) provides a natural
compilation target: each certified rewrite becomes a runtime
primitive with machine-checked correctness.

\section{Process-Algebraic and Multi-Agent Extensions}

The typed world-model encoding currently covers single-agent
reasoning.  Extending to multi-agent settings---where agents
communicate partial evidence, negotiate under uncertainty, and
maintain independent world models that must be fused---requires
process-algebraic infrastructure: confluence modulo associativity and
commutativity, communication protocols with evidence-typed channels,
and compositional verification of multi-agent inference systems.

\section{Broader Benchmark Coverage}

The present empirical validation covers specific domains.
Drug-repurposing knowledge-graph benchmarks (Hetionet/Rephetio),
genomic regulatory pipelines, large-scale knowledge-base reasoning,
and real-time sensor fusion at scale remain as identified but
unfinished empirical targets.  Each connects to a specific
theoretical capability of the calculus and would test the practical
reach of the formal guarantees.
